\section*{अध्याय \ref{ch:b3:stem-cells} --- मूल कोशिकाएँ, पुनर्जनन और वृद्धावस्था}

\begin{solution}{exo:b3:stem-cells:1}
\omterm{def:b3:stem-cells:stem}{स्व-नवीकरण}: कोई ऐसा विभाजन जो कम से कम एक ऐसी संतति बनाए जो अब भी \omterm{def:b3:stem-cells:stem}{मूल
कोशिका} हो --- यानी क्रिप्ट के आधार की \emph{Lgr5} कोशिकाएँ। शक्ति: उन
भवितव्यों की परास जो कोई कोशिका अब भी अपना सकती है --- क्रिप्ट की \omterm{def:b3:stem-cells:stem}{मूल
कोशिका} बहुशक्त है, जो अवशोषी, चषक, आंत्र-अंतःस्रावी, पैनेथ तथा गुच्छ
कोशिकाएँ बनाती है। \omterm{def:b3:stem-cells:stem}{कोटर}: \omterm{prop:b3:stem-cells:crypt}{पैनेथ कोशिकाएँ} और उनके Wnt, Notch तथा EGF
संकेत, जो मूलत्व को अपनी जगह थामे रखते हैं। \omterm{def:b3:stem-cells:stem}{संक्रमण-प्रवर्धी कोशिका}:
कोई प्रतिबद्ध संतति, जो रसांकुर पर ऊपर जाते हुए विभेदित होने से पहले उस
क्रिप्ट में चार-पाँच बार और बँटती है।
\end{solution}

\begin{solution}{exo:b3:stem-cells:2}
एकल मज्जा-कोशिकाओं ने ऐसी प्लीहा-बस्तियाँ बनाईं जिनमें लाल कोशिकाएँ,
कणिकाणु तथा बिंबाणु थे: यानी एक कोशिका, कई वंश-परंपराएँ --- यानी
बहुशक्तता, और उस रैखिक मात्रा-अनुक्रिया ने दिखाया कि हर बस्ती किसी एक
ही कोशिका से आई। \omterm{def:b3:stem-cells:stem}{स्व-नवीकरण} के लिए कोई दूसरा चूहा चाहिए था: किसी बस्ती
की कोशिकाएँ, किसी दूसरे विकिरणित ग्रहीता में अंतःक्षिप्त करने पर, नई
बस्तियाँ बना देती थीं, इसलिए उस संस्थापक कोशिका ने अपनी ही क्षमता वाली
संततियाँ बनाई थीं।
\end{solution}

\begin{solution}{exo:b3:stem-cells:3}
\emph{Oct4}, \emph{Sox2}, \emph{Klf4}, \emph{c-Myc}। उस तंतुकोरक के
\omterm{def:b3:chromatin-epigenetics:nucleosome}{क्रोमैटिन} ने \omterm{def:b3:stem-cells:pluripotency}{सर्वशक्तता} के जीन \omterm{def:b3:chromatin-epigenetics:nucleosome}{विषमक्रोमैटिन} तथा DNA-मेथिलन के पीछे बंद
कर रखे हैं, इसलिए उन कारकों को पहले कुछ ही पहुँच-योग्य स्थल मिलते हैं
और उन्हें अनेक \omterm{def:b3:cell-cycle-apoptosis:cdk}{कोशिका-चक्रों} भर पुनर्रचक एंज़ाइम भर्ती करनी पड़ती हैं;
उस कोशिका को अपना ही कार्यक्रम भी बंद करना पड़ता है, जरण से बचना पड़ता
है, और किसी ऐसी प्रसंभाव्य मध्यवर्ती अवस्था से गुज़रना पड़ता है जिसमें
अधिकांश कोशिकाएँ अटक या मर जाती हैं। केवल वही दुर्लभ कोशिका, जिसमें हर
क़दम सफल हो, हफ़्तों बाद सर्वशक्त निकलती है।
\end{solution}

\begin{solution}{exo:b3:stem-cells:4}
प्लेनेरिया सर्वशक्त वयस्क \omterm{def:b3:stem-cells:stem}{मूल कोशिकाएँ}, यानी \omterm{def:b3:stem-cells:regeneration}{नियोब्लास्ट}, रखता है, जो
उस घाव तक प्रवास करती हैं और जो भी ग़ायब हो वह बना देती हैं।
सैलामैंडर के पास ऐसा कोई पूल नहीं: उस ठूँठ की विभेदित कोशिकाएँ
विभेदन-निरसित होकर \omterm{def:b3:stem-cells:stem}{पूर्वगामी} बनती हैं, जो निवासी ऊतक-मूल कोशिकाओं के साथ
मिलकर कोई ऐसा \omterm{def:b3:stem-cells:regeneration}{ब्लास्टीमा} बनाती हैं जो उपांग-कार्यक्रम फिर चला देता है।
क्रागल के एक-एक अंकित ऊतक के कलमों ने दिखाया कि \omterm{def:b3:stem-cells:regeneration}{ब्लास्टीमा} की कोशिकाएँ
अपने उद्गम पर सच्ची रहती हैं --- पेशी पेशी बनाती है, उपास्थि उपास्थि
--- इसलिए वह \omterm{def:b3:stem-cells:regeneration}{ब्लास्टीमा} वंश-सीमित \omterm{def:b3:stem-cells:stem}{पूर्वगामियों} का कोई मिश्रण है, कोई
सर्वशक्त पिंड नहीं।
\end{solution}

\begin{solution}{exo:b3:stem-cells:5}
$(11 - 5)/0.075 = 80$ द्विगुणन। \qty{8}{kb} से: $(8 - 5)/0.075 = 40$।
\qty{60}{\%} भरपाई के साथ शुद्ध हानि \qty{30}{bp} है: $6000/30
= 200$।
\end{solution}

\begin{solution}{exo:b3:stem-cells:6}
चौदह मूल-कोशिका विभाजन दिन भर में 28 संततियाँ देते हैं; 14 को \omterm{def:b3:stem-cells:stem}{मूल
कोशिकाएँ} बने रहना है, इसलिए 14 प्रतिबद्ध होती हैं --- यानी प्रति \omterm{def:b3:stem-cells:stem}{मूल
कोशिका} प्रति दिन एक प्रतिबद्ध संतति। हर एक $2^{4} = 16$ कोशिकाओं तक
प्रतिबद्ध होती है: $14\times 16 = 224$ रसांकुर-कोशिकाएँ प्रति दिन, यानी
देखी गई \num{300} की कोटि (कोई पाँचवाँ संक्रमण-विभाजन \num{448} देता
है)। मूल-कोशिका तह चौदह कोशिकाओं जितने विभाजन देती है, और संक्रमण-तह
बाकी \num{200}।
\end{solution}

\begin{solution}{exo:b3:stem-cells:7}
दिन भर में $2\times 10^{11}/2^{20} = 1.9\times 10^{5}$ प्रतिबद्ध
संततियाँ; \num{10000} \omterm{def:b3:stem-cells:stem}{मूल कोशिकाओं} के साथ, प्रति \omterm{def:b3:stem-cells:stem}{मूल कोशिका} प्रति दिन
19 --- जबकि विभाजन साल में एक। बीस विभाजनों की वह गिनती \omterm{def:b3:stem-cells:stem}{मूल कोशिका} से
लाल कोशिका तक प्रति वंश-परंपरा है, पर वे मध्यवर्ती \omterm{def:b3:stem-cells:stem}{पूर्वगामी} एक-बार-के
नहीं हैं: बहुशक्त तथा प्रतिबद्ध \omterm{def:b3:stem-cells:stem}{पूर्वगामी} पूल अपने ही विभाजनों से
हफ़्तों तक अपने को बनाए रखते हैं, इसलिए लगभग सारा विभाजन उन्हीं
\omterm{def:b3:stem-cells:stem}{पूर्वगामी} तहों में होता है और उस \omterm{def:b3:stem-cells:stem}{मूल कोशिका} को बिरले ही बुलाया जाता है।
\end{solution}

\begin{solution}{exo:b3:stem-cells:8}
$\mu(50) = 10^{-3}\mathrm{e}^{1.74} = 5.7\times 10^{-3}$; $\mu(70) =
10^{-3}\mathrm{e}^{3.48} = 0.032$; $\mu(90) = 10^{-3}\mathrm{e}^{5.22} =
0.19$ प्रति वर्ष। $\mu_{0}/\gamma = 0.0115$ के साथ: $S(70) = \exp[-0.0115
\times 31.5] = 0.70$, $S(90) = \exp[-0.0115\times 184] = 0.12$। मध्यक:
$30 + \ln(61)/0.087 = 77$। $\mu_{0}$ आधा करने पर: $\ln(121)/0.087 = 55$,
मध्यक 85 --- यानी एक द्विगुणन-काल, आठ वर्ष, की बढ़त। $\gamma$ आधा करके
$0.0435$ करने पर: $\ln(31)/0.0435 = 79$, मध्यक 109: उस घातांक को धीमा
करना उस अचर को आधा करने से चार गुना अधिक देता है।
\end{solution}

\begin{solution}{exo:b3:stem-cells:9}
बचे तिहाई को तिगुना होना है: प्रति कोशिका $\log_{2}3 = 1.6$ विभाजन।
\omterm{def:b3:stem-cells:regeneration}{प्रतिपूरक वृद्धि}, क्योंकि मौजूद पालियाँ अपनी जगह कोशिकाओं के विभाजन से
बड़ी हो जाती हैं; और हटाई गई पालियाँ, अपनी नलिकाओं, वाहिकाओं तथा वास्तु
समेत, फिर नहीं बनतीं --- वह यकृत अपना द्रव्यमान तथा प्रकार्य लौटा लेता
है, अपना आकार नहीं।
\end{solution}

\begin{solution}{exo:b3:stem-cells:10}
हर घटना पर उस क्लोन का आकार $k$ बराबर प्रायिकता से $k + 1$ या $k - 1$
हो जाता है, इसलिए उसका प्रत्याशित आकार कभी नहीं बदलता; और जब वह चलन
समाप्त होता है तब वह $N$ (क़ब्ज़ा) या $0$ (लोप) होता है, और प्रत्याशा
$N\,P + 0 = k$ से $P = k/N$ मिलता है। कोई एक कोशिका: $1/14$। किसी एक \omterm{def:b3:stem-cells:stem}{मूल
कोशिका} में उठा कोई उदासीन उत्परिवर्तन अपने क्रिप्ट में $1/14$ की
प्रायिकता से स्थिर होता है और अन्यथा महीनों में खो जाता है --- अपवाह
चौदह में से तेरह उत्परिवर्तन साफ़ कर देता है।
\end{solution}

\begin{solution}{exo:b3:stem-cells:11}
वह चलन अब पक्षपाती है: वह क्लोन जितनी बार कोई जगह खोता है उससे अधिक बार
पाता है, इसलिए उसका प्रत्याशित आकार बढ़ता है और $k$ के बढ़ने पर क़ब्ज़े
की प्रायिकता निश्चितता की ओर चढ़ती जाती है --- चौदह के किसी \omterm{def:b3:stem-cells:stem}{कोटर} में
\qty{10}{\%} लाभ के लिए किसी एक कोशिका के $1/14$ से मोटे तौर पर तिगुनी,
और क्लोन के कुछ जगहें पा लेने पर उससे भी ऊँची। दशकों में ऐसे
उत्परिवर्तनों वाली \omterm{def:b3:stem-cells:stem}{मूल कोशिकाएँ} (मज्जा में \emph{DNMT3A}, \emph{TET2})
अपने \omterm{def:b3:stem-cells:stem}{कोटरों} पर क़ब्ज़ा कर लेती हैं, और सत्तर तक पाँचवें भाग लोगों का
रक्त बड़े हिस्से में किसी एक क्लोन से आता है। वह \omterm{def:b3:cancer-biology:hallmarks}{कर्कट} नहीं है: वे
कोशिकाएँ अब भी विभेदित होती हैं और अपनी पदानुक्रमी मानती हैं। पर अब
अगले उत्परिवर्तन के उठने के लिए अरबों का कोई क्लोन मौजूद है, और
रक्तकर्कट का ख़तरा दस गुना बढ़ जाता है।
\end{solution}

\begin{solution}{exo:b3:stem-cells:12}
मध्यक $= 30 + \gamma^{-1}\ln(1 + \gamma\ln 2/\mu_{0})$: $\mu_{0} =
10^{-2}$ के लिए, $30 + \ln 7/0.087 = 52$; $10^{-3}$: 77; $10^{-4}$: $30 +
\ln 604/0.087 = 104$। शुद्ध गोंपर्ट्ज़ रूप में हर दस गुनी कटौती वही
$\ln 10/\gamma = 26$ वर्ष जोड़ती है। व्यवहार में प्रतिफल इसलिए घटता है
कि ऐतिहासिक बढ़तें किसी आयु-निरपेक्ष पद को हटाने से आई थीं ---
संक्रमण, प्रसव, दुर्घटना, यानी मेकहम का अचर $A$, जो $\mu = A
+ \mu_{0}\mathrm{e}^{\gamma t}$ में है --- और जब $A$ उन
आयुओं पर, जो मायने रखती हैं, उस चरघातांक के सामने छोटा पड़ जाए, तब उसे
और घटाने से लगभग कुछ नहीं बदलता; और वह आंतरिक $\mu_{0}$ कहीं कम खिसका
है। $\gamma$ में \qty{20}{\%} की कटौती ($0.070$ तक) देती है
$30 + \ln 49/0.070 = 86$: यानी नौ वर्ष, जो मृत्यु के कुछ कारण टालकर
नहीं बल्कि हर आयु पर मिलते हैं।
\end{solution}

\begin{solution}{pb:b3:stem-cells:1}
\textbf{1.} दिन भर में $3.5\times 10^{11}$, प्रति सेकंड $4\times 10^{6}$।
\textbf{2.} प्रति प्रतिबद्ध संतति $2^{20} \approx 10^{6}$ कोशिकाएँ;
दिन भर में $3.5\times 10^{11}/10^{6} = 3.3\times 10^{5}$ प्रतिबद्ध संततियाँ।
\textbf{3.} प्रति \omterm{def:b3:stem-cells:stem}{मूल कोशिका} प्रति दिन $33$, प्रति वर्ष \num{12000}।
\textbf{4.} उन \omterm{def:b3:stem-cells:stem}{पूर्वगामी} पूलों को अपने को बनाए रखना पड़ता है: बहुशक्त
तथा प्रतिबद्ध \omterm{def:b3:stem-cells:stem}{पूर्वगामी} हफ़्तों तक बँटते हैं और अपनी ही संख्याएँ थामे
रहते हैं, इसलिए उन बीस में से लगभग हर विभाजन उन तहों में होता है जो
अपना नवीकरण करती हैं, और वह \omterm{def:b3:stem-cells:stem}{मूल कोशिका} बिरले ही कोई प्रतिबद्ध संतति
देती है --- कोई साल में एक बार की कोटि पर।
\textbf{5.} $10 - 20\times 0.06 = \qty{8.8}{kb}$: यानी
$L_{\text{क्रांतिक}}$ से कहीं ऊपर, और वह लाल कोशिका फिर कभी बँटती ही नहीं;
वह भंडार केवल उन \omterm{def:b3:stem-cells:stem}{पूर्वगामियों} के लिए मायने रखता है जो फिर काम में लिए
जाते हैं।
\textbf{6.} शुद्ध हानि $0.3\times 60 = \qty{18}{bp}$; $6000/18 = 333$
विभाजन; साल में एक पर, $333$ वर्ष --- यानी किसी जीवन से कहीं आगे,
इसलिए टीलोमियर उस \omterm{def:b3:stem-cells:stem}{मूल कोशिका} को स्वयं सीमित नहीं करते; दूसरी क्षति
करती है।
\textbf{7.} दिन भर में $14$ विभाजन; $7$ खोई \omterm{def:b3:stem-cells:stem}{मूल कोशिकाओं} की जगह भरते हैं
और $7$ प्रतिबद्ध होते हैं; $7\times 2^{4} = 112$ रसांकुर-कोशिकाएँ प्रति
दिन।
\textbf{8.} दिन में एक पर प्रति कोशिका $N^{2} = 196$ घटनाएँ: यानी लगभग
\qty{200}{d}, छह-सात महीने --- यानी वही समय जिसमें कन्फ़ेटी क्रिप्ट
एकरंगी हुए।
\textbf{9.} $1/14 \approx \qty{7}{\%}$; $10^{7}/14 \approx 7\times 10^{5}$
क्रिप्ट।
\textbf{10.} अपवाह हर चौदह में से तेरह उत्परिवर्तन महीनों में बाहर फेंक
देता है। जीवन भर का \omterm{def:b3:stem-cells:stem}{असममित विभाजन} हर उत्परिवर्ती \omterm{def:b3:stem-cells:stem}{मूल कोशिका} को जीवन भर
रखता, इसलिए उत्परिवर्तन हर वंश-परंपरा में बिना कभी साफ़ हुए जमा होते
जाते।
\textbf{11.} मूलत्व उस \omterm{def:b3:stem-cells:stem}{कोटर} में कोई स्थिति और संकेतों का कोई समुच्चय
है, किसी कोशिका का कोई स्थायी गुण नहीं: जो भी कोशिका पैनेथ संकेतों तक
पहुँचे वह वह भूमिका ले सकती है।
\textbf{12.} किसी एक क्षण उन \omterm{def:b3:stem-cells:stem}{मूल कोशिकाओं} में कोई बहुरंगी अंक प्रेरित
कीजिए और उन क्रिप्टों का पीछा कीजिए। सममित अपवाह: क्रिप्ट महीनों में
एकरंगी हो जाते हैं। \omterm{def:b3:stem-cells:stem}{असममित विभाजन}: हर क्रिप्ट अपने सारे रंग अनिश्चित काल
तक रखता है, हर एक रसांकुर पर चढ़ती किसी पतली फीते की तरह।
\textbf{13.} $6000/60 = 100$ द्विगुणन।
\textbf{14.} लगभग $75$ बचे: वह कोशिका विभाजन गिनती है, पंचांग का समय
नहीं, क्योंकि जमाव में बीते दशक की कोई क़ीमत नहीं पड़ी।
\textbf{15.} बीस के बाद $(8000 - 4000)/30 = 133$ वर्ष --- यानी आयु 153।
माध्य श्वेताणु टीलोमियर-लंबाई स्वयं जीवन सीमित नहीं करती; सबसे छोटे
टीलोमियर, और वे दूसरे लक्षण, अधिक मायने रखते हैं।
\textbf{16.} \omterm{def:b3:cancer-biology:telomerase}{टीलोमरेज़} एक अवरोध, यानी \omterm{thm:b3:stem-cells:hayflick}{प्रतिकृतिक जरण}, हटा देती है, और
इसीलिए वह लगभग हर \omterm{def:b3:cancer-biology:hallmarks}{कर्कट} में मिलती है, जिसे असीमित विभाजन चाहिए; पर
जिन तंतुकोरकों को \omterm{def:b3:cancer-biology:telomerase}{टीलोमरेज़} दी गई उन्होंने अपने जाँच-बिंदु, संपर्क-निरोध
तथा अक्षत p53 रखे, इसलिए अकेले असीमित विभाजन ने उन्हें \omterm{def:b3:cancer-biology:hallmarks}{कर्कट} नहीं
बनाया। आवश्यक, पर्याप्त नहीं।
\textbf{17.} $\mu(40) = 2.4\times 10^{-3}$, $\mu(60) = 0.014$, $\mu(80) =
0.077$, $\mu(100) = 0.44$ प्रति वर्ष। $S = \exp[-0.0115(\mathrm{e}^{\gamma
(t-30)} - 1)]$: 40 पर $0.98$, 60 पर $0.87$, 80 पर $0.42$, और 100 पर
$0.006$।
\textbf{18.} मध्यक $30 + \ln 61/0.087 = 77$। दस प्रतिशत तब जीवित रहते
हैं जब $\mathrm{e}^{\gamma(t-30)} = 1 + \gamma\ln 10/\mu_{0} = 201$: $t = 30 +
5.3/0.087 = 91$।
\textbf{19.} \qty{0.67}{mm/d}; 60 दिनों में $\log_{2}1000 = 10$
द्विगुणन, यानी हर छह दिन में एक।
\textbf{20.} लगभग $10^{8}$ विभाजन, इसलिए प्रति \omterm{def:b3:stem-cells:regeneration}{पुनर्जनन} $0.1$ अर्बुदजनक
उत्परिवर्तन --- यानी दस उपांगों में एक। ऐक्सोलोटल में \omterm{def:b3:cancer-biology:hallmarks}{अर्बुद} लगभग कभी
नहीं पनपते: वे ठंडे-रक्ती और धीमे हैं, उनका अर्बुद-दमन मज़बूत है, और
उस \omterm{def:b3:stem-cells:regeneration}{ब्लास्टीमा} की कोशिकाएँ वंश-सीमित रहती हैं तथा बँटती रहने के बजाय फिर
से विभेदित हो जाती हैं; \omterm{def:b3:stem-cells:regeneration}{पुनर्जनन} करता ऊतक तो प्रेरित \omterm{def:b3:cancer-biology:hallmarks}{अर्बुद} भी दबा देता
है। किसी स्तनी के पास अधिक कोशिकाएँ, ऊँचा तापमान तथा उपापचय-दर, और कोई
लंबा जीवन है जिसमें कोई बच निकला क्लोन बढ़ सकता है।
\textbf{21.} उस ठूँठ की तंत्रिका काटकर वहाँ उम्मीदवार प्रोटीन छोड़ता कोई
मनका बैठाइए, या तंत्रिका-अनुकूलित ऊतक कलमित कीजिए: \omterm{def:b3:stem-cells:regeneration}{पुनर्जनन} लौट आए तो
वह कोई विसरणशील कारक है। न्यूटों तथा ऐक्सोलोटलों में पहचाने गए
उम्मीदवार: अग्र-प्रवणता प्रोटीन nAG, न्यूरेगुलिन-1, और FGF।
\textbf{22.} \omterm{prop:b3:stem-cells:crypt}{Wnt संकेतन} पूँछ पर ऊँचा और सिर पर नीचा है, और किसी टुकड़े
के अग्र सिरे का घाव सामान्यतः सिर इसलिए बनाता है कि वहाँ Wnt नीचा होता
है; किसी Wnt निरोधक को घटाने से Wnt हर जगह चढ़ जाता है, दोनों घाव
``पश्च'' पढ़ते हैं, और दो पूँछें उग आती हैं। वह टुकड़ा अपना अक्ष अपनी
पेशी में बची प्रवणता से जानता है, और वह घाव उस प्रवणता को उसी के सापेक्ष
फिर खड़ा कर देता है।
\textbf{23.} बचे तिहाई को तिगुना होना है: हर यकृतकोशिका एक बार बँटती है
(दो-तिहाई तक) और आधी फिर बँटती हैं --- यानी माध्य $1.6$ विभाजन। बची
पालियाँ फूलकर मूल द्रव्यमान तक आ जाती हैं; और ग़ायब पालियाँ तथा उनकी
नलिकाएँ और वाहिकाएँ फिर नहीं बनतीं।
\textbf{24.} जो \omterm{def:b3:stem-cells:regeneration}{पुनर्जनन} घिसे तथा क्षतिग्रस्त ऊतक की जगह भर दे वह उस
जमाव को धीमा कर देता जिसे वह घातांक नापता है, यानी $\gamma$ को गिरा
देता और उस वक्र को चपटा कर देता --- और ऐक्सोलोटल जरण बहुत कम दिखाते
हैं। उसकी क़ीमत किसी गरम, बड़े, दीर्घजीवी शरीर में \omterm{def:b3:cancer-biology:hallmarks}{कर्कटों} से आता कोई
ऊँचा अचर पद होती।
\textbf{25.} हेफ़्लिक गिनती लगभग $100$ द्विगुणन; किसी
टीलोमरेज़-सहित \omterm{def:b3:stem-cells:stem}{मूल कोशिका} के पास कोई $330$ विभाजन हैं, यानी साल में एक
पर तीन सदियों से अधिक; और मध्यक जीवनकाल $77$ वर्ष।
\end{solution}
